\chapter{记忆语义模型}

\section{记忆域的职责}

记忆域保存主体决定跨会话保留的经历与认知，并拥有与这些认知直接相关的 Episode、Journal、CognitiveSchema、Tag 和证据/修订合同。聊天逐字稿、原始材料、观察事件和外部知识仍由各自对象及来源管理；主体接触材料后可以形成长期认知，也可以只把材料用于当前运行时。\cite{nous-research-package}

目标形成流程为：

\begin{center}
\small
来源快照 $\rightarrow$ 当前相关认知 $\rightarrow$ 提案 $\rightarrow$ 主张/证据/语义指向映射 $\rightarrow$ 所有者核验 $\rightarrow$ 对象纪元与当前版本核验 $\rightarrow$ 不可变修订
\end{center}

模型生成默认只形成提案；只有特定所有者合同明确允许时，确定性流程才可直接提交。每个记忆修订保留证据引用（EvidenceRef）、生产者签名（ProducerSignature）和来源谱系。

\section{材料、接触记录、派生表示与证据}

\begin{description}
\item[原始材料（Artifact）] 稳定的原始或生成内容；大文件可保存在不可变对象库中。
\item[观察事件（ObservationOccurrence）] 主体在特定时间实际接触材料或事件的一次记录。相同字节内容可以对应多次观察。
\item[派生表示（DerivedRepresentation）] 由光学字符识别、解析器或模型生成的版本化解释，保存生成者、版本和来源区域映射。
\item[证据引用（EvidenceRef）] 指向观察事件、来源区域、派生表示或派生区域的来源句柄。
\item[记忆对象（Memory）] 主体接受并决定长期保留的认知对象，具有稳定对象身份和不可变修订。
\end{description}

\begin{ReportWideFigure}
\centering
\begin{tikzpicture}[
  node distance=0.5cm,
  block/.style={draw=NWBlue,rounded corners,align=center,text width=3.2cm,minimum height=1.05cm,inner sep=5pt},
  small/.style={draw=NWRule,rounded corners,align=center,text width=3.2cm,minimum height=0.9cm,inner sep=4pt},
  arrow/.style={-{Stealth[length=2mm]},thick,color=NWBlue}
]
\node[block] (artifact) {原始材料\\稳定来源内容};
\node[block,right=of artifact] (encounter) {观察事件\\主体的一次接触};
\node[block,right=of encounter] (derived) {派生表示\\版本化解释};
\node[small,below=1.1cm of encounter] (evidence) {证据引用\\来源谱系};
\node[block,right=of evidence] (proposal) {认知提案\\主张与语义指向};
\node[block,right=of proposal] (revision) {领域所有者核验\\不可变修订};
\draw[arrow] (artifact) -- (encounter);
\draw[arrow] (encounter) -- (derived);
\draw[arrow] (encounter.south) -- (evidence.north);
\draw[arrow] (derived.south) |- (evidence.east);
\draw[arrow] (evidence) -- (proposal);
\draw[arrow] (proposal) -- (revision);
\end{tikzpicture}
\captionof{figure}{材料进入主体认知的边界。派生表示不替代观察事件，也不会自动成为记忆。}
\label{fig:memory-formation}
\end{ReportWideFigure}

认知对象由自身身份与修订标识；证据引用说明支撑该对象的来源。多个派生文本若共享同一来源祖先，证据独立性仍按共同来源谱系计算。检索保留来源链，区分独立佐证与同源转述。

\section{身份、修订与生命周期}

对同一认知对象的修正、澄清或重释，通常创建该对象的新修订；独立的新经历创建新对象；矛盾认知通常以显式矛盾关系保留，不通过把一方标记为“已取代”来隐藏。修订、当前版本、对象纪元和生命周期分别承担内容版本、当前指针、并发栅栏和生命周期状态。\cite{nous-research-package}

对象级状态与内容修订保持分离。抑制、恢复、固定访问策略或其他不改变认知内容的操作可以推进对象纪元而不伪造新的语义修订；真正内容变化才创建新 revision。

\section{正交记忆分类}

目标设计不再使用互斥的 \texttt{Specific / Integrative / Procedural} 作为单一 \texttt{MemoryClass}。它们混合了“如何形成”和“认知内容承担什么作用”两个维度。基础合同改为至少两个正交轴：

\begin{description}
\item[认知角色（cognitive role）] \texttt{Experiential} 表达主体经历或遇到过什么；\texttt{Declarative} 表达主体长期保留的事实、命题或解释性认知；\texttt{ProceduralExperience} 表达从经验中记住的做法、步骤或策略性经验。ProceduralExperience 仍是认知，不是可执行行为策略。
\item[形成方式（formation mode）] \texttt{Grounded} 主要由一个有界经历/来源上下文形成，可带额外佐证；\texttt{Synthesized} 跨多个认知对象或独立来源综合形成。
\end{description}

认知角色更接近对象的稳定身份性质；形成方式属于形成该 revision 的来源语义。一个 Synthesized ProceduralExperience 因而是合法组合，不需要强迫归入互斥类别。“Specific”不再作为全局对象类别；其具体程度由 evidence、Episode、适用范围和内容本身表达。

暂不增加通用 \texttt{Metacognitive} 记忆类别。关于主体自身的长期稳定理解可能属于 Self；只有出现明确独立生命周期与查询合同后，再决定是否扩充 Memory cognitive role。

\section{时间语义}

至少需要区分：

\begin{itemize}
\item 事件发生时间（occurred\allowbreak\_time）；
\item 主体观察时间（observed\allowbreak\_at）；
\item 主张有效时间（valid\allowbreak\_time）；
\item 系统记录时间（recorded\allowbreak\_at）；
\item 主体权威提交序号（authority\allowbreak\_seq）；
\item 认知形成时间（formed\allowbreak\_at）。
\end{itemize}

记忆的观察时间沿证据谱系追溯到真实接触事件。读取视图可展示派生的最早和最晚时间；范围值保留为读取结果，不写作新的权威事实。未知时间保持未知，并由查询策略显式处理。

\section{语义指向、认识状态与来源提及}

认识状态、证据支持角色、来源谱系、认知角色和生命周期分别建模。材料提及（EntityMention）只记录来源中真实出现的实体；MemoryRevision aboutness 独立记录该认知涉及哪些 EntityRef 及其角色。认知指向不能附着到任意一条 evidence 上，也不能由第一条证据或向量相似度猜测。

当前代码已经将记忆修订到实体的语义关联与来源 mention 分开，并把 contradiction 从 revision lifecycle 中拆出；这些语义应直接继承到下一版合同。当前代码仍保留互斥 \texttt{MemoryClass}，属于目标实现需要替换的部分。\cite{nous-code-snapshot}

\section{可访问性、遗忘与管理生命周期}

持久认知是否存在、当前是否允许普通访问、当前是否容易被检索、以及与当前查询是否相关是四个不同问题：

\begin{center}
Authority/Durability $\neq$ Accessibility $\neq$ Relevance $\neq$ Suppression/Purge
\end{center}

Authority 保存认知内容及显式管理意图；\texttt{AUTO} accessibility 应在查询时根据时间、类型化使用、线索和预算产生可解释的 AccessibilityDecision，而不是把一个随时间变化的“记忆强度真值”永久写入认知对象。固定（Pinned）或精确访问限制等显式管理意图可以持久保存。

当前实现中的 90/365 天 NORMAL/DEEP/EXPLICIT 阈值是可运行实验基线，不是认知规律。Memory Reference Profile 不以硬时间阈值作为正确性前提：AUTO 默认不因年龄永久排除普通合法记忆；时间和 meaningful use 可以作为软先验、深度预算或诊断输入。后续通过无衰减、仅时间、时间加类型化使用及其他候选模型比较后再冻结默认策略。

抑制（suppression）阻止普通认知消费但保留权威内容；purge 是显式管理权限驱动的跨 Authority/Serving/cache/material 流程；runtime eviction 只移出当前工作集。自然语言“忘掉”本身不具有 destructive authority。